\subsection{A family of ordinary genus-two curves in characteristic two}

\begin{claim}{A.C2.1}{Proved}
Let $K$ be finitely generated over $\F_2$, let $a,b,c\in K^\times$,
and let $C_{a,b,c}$ be the smooth proper curve with function field
\[
 y^2+y=a x+\frac b{x}+\frac c{x+1}.
\]
It is an ordinary genus-two curve. If
\[
                    (a+b+c)^3+abc=0,                         \tag{C2.1}
\]
then its geometrically prime-to-two point-section map is injective.
The same conclusion holds for every hyperbolic open subcurve and after
any finite extension of $K$. It also holds for any curve that becomes
isomorphic to such a curve after a finite extension.

Over an algebraic closure, condition (C2.1) is equivalent to
\[
 b=a v^3,\qquad c=a(v+1)^3\qquad(v\ne0,1).                    \tag{C2.2}
\]
For each such $v$ there is a connected cyclic \'etale cover $D\to C$
of degree three and two maps $D\to E$ to one supersingular elliptic
curve whose product is injective on geometric points.
\end{claim}
\begin{proof}
Work initially over an algebraic closure $\Omega$ of $K$. The three
simple Artin--Schreier poles are $0,1,\infty$. Their different
exponents are two, so Riemann--Hurwitz gives genus two. Write
\[
 h=x^2+x,\qquad Y=h y,\qquad
 f=a x^5+(a+b+c)x^3+c x^2+b x.
\]
The equation is $Y^2+hY=f$. There is one point $P_\infty$ above
infinity, with pole orders two and five for $x$ and $Y$, respectively.
The regular differentials $dx/h$ and $x\,dx/h$ form a basis. Cartier
fixes each of them: write $1/h=h/h^2$ and use
$\mathcal C((x^2+x)dx)=dx$ and
$\mathcal C((x^3+x^2)dx)=x\,dx$ in characteristic two.
Thus Cartier is bijective and the curve is ordinary.

For the parameterization put $B=b/a$, $C=c/a$. Equation (C2.1)
becomes $(1+B+C)^3+BC=0$. If $B+C\ne0$, set
$v=(C+1)/(B+C)$. Direct substitution gives
$v^2+v=B+C+1$ and $v^3=B$, whence $(v+1)^3=C$.
If $B=C$, the equation implies $B=C=1$; either root of
$v^2+v+1=0$ works. Conversely (C2.2) gives (C2.1) by expansion.
Nonvanishing of $b,c$ excludes $v=0,1$.

Choose $d\in\Omega$ satisfying $d^2+d=a^2v^2+a$, and set
\[
 u=x^2+v^2,\qquad
 A=a x^3+d x^2+(d+a+b+c)x+b.
\]
A coefficient calculation gives
\[
                  A^2+hA+f=a^2u^3.                          \tag{C2.3}
\]
Let $F=(A+Y)/a$. The hyperelliptic involution $\iota$ sends
$Y$ to $Y+h$, so $F\iota(F)=u^3$. At $x=v$, the map
$C\to\mathbf P^1$ is unramified, since $h(v)=v(v+1)\ne0$.
Of its two points, precisely one, denoted $R$, has $F(R)=0$.
The conjugate value is $h(v)/a\ne0$. Because $u$ has order two
there, (C2.3) gives
\[
                 \operatorname{div}(F)=6R-6P_\infty.
\]
The normalization in $z^3=F$ is therefore \'etale. It is connected:
if $F$ were a cube, its cube root would have its only pole at
$P_\infty$ of order two and would lie in
$L(2P_\infty)=\langle1,x\rangle$. Its cube could not contain the
nonzero $Y$ coefficient of $F$. Denote this cyclic cover by $D$;
it has genus four.

Put $w=z+u/z$. The involution $z\mapsto u/z$ gives a map of
degree two from $D$ to
\[
 E:\quad a\bigl(W^3+(X^2+v^2Z^2)W\bigr)=X^2Z+XZ^2,
 \qquad O=[1:0:0].                                          \tag{C2.4}
\]
The derivatives with respect to $X$ and $Z$ are $Z^2$ and
$X^2$. A singular projective point would therefore have $X=Z=0$,
and the derivative with respect to $W$ would then force $W=0$.
Thus this is a smooth plane cubic.

We verify supersingularity without a characteristic restriction beyond
two. For a smooth plane cubic $G(X,W,Z)=0$ in characteristic two,
the scalar of Frobenius on $H^1(E,\mathcal O_E)$ is the coefficient
of $XWZ$ in $G$, in the usual one-dimensional Cech description.
Indeed the exact sequence
$0\to\mathcal O_{\mathbf P^2}(-3)\xrightarrow{G}
\mathcal O_{\mathbf P^2}\to\mathcal O_E\to0$
identifies this group with $H^2(\mathcal O_{\mathbf P^2}(-3))$,
whose generator is represented by $1/(XWZ)$. Frobenius takes it
to $G/(X^2W^2Z^2)$; every cubic monomial except $XWZ$ gives a
zero class in this Cech group. In (C2.4) that coefficient is zero.
Hence the elliptic curve has $2$-rank zero, and
$E(\Omega)[2^\infty]=0$.

Let $\zeta\ne1$ satisfy $\zeta^3=1$ and put $q=u/z$. The two maps
\[
 r_0=(x,z+q),\qquad r_1=(x,\zeta z+\zeta^{-1}q):D\to E
\]
extend to the smooth proper curves. We check their joint injectivity.
For finite $x\ne v$, the two image coordinates recover $z,q$, since
\[
 \det\begin{pmatrix}1&1\\ \zeta&\zeta^{-1}\end{pmatrix}=1.
\]
Together with $Y=a z^3-A$ they recover the point. At $x=v$, let
$\mathcal R=\{s:s^3=h(v)/a\}$; it consists of three distinct
nonzero elements. The three points above $R$ have $z=0$, $q\in
\mathcal R$, and their two $w$-coordinates are
$(q,\zeta^{-1}q)$. The other three have $q=0$, $z\in\mathcal R$,
and coordinates $(z,\zeta z)$. These six pairs are distinct.
At infinity the three residues of $z/x$ are $1,\zeta,\zeta^2$,
while the residues of $q/x$ are their inverses. Writing
$T=[1:1:0]\in E$, their image pairs are
\[
                         (O,T),\quad(T,T),\quad(T,O).
\]
Every finite point maps to the affine chart $Z=1$, so no other
collision with these pairs occurs. The two maps are jointly injective
on all geometric points.

All choices are defined after a finite extension of $K$. A hypothetical
collision of point sections persists over that extension and over a
further extension making a chosen lift to $D$ rational. The fixed-fibre
argument of A.3 supplies a matching lift whose Jacobian difference is
$2$-primary by A.1. Applying $r_0$ and $r_1$ kills this difference,
because $E(\Omega)[2^\infty]=0$. Joint injectivity forces the two
lifts, and then the original points, to coincide. This is the argument
of A.10, with the same Kummer input valid over finitely generated
fields. A collision on an open subcurve maps to one on its proper
compactification. These observations also prove the assertions about
finite extensions and forms.
\end{proof}

\begin{claim}{A.C2.2}{Proved}
For the ordinary genus-two curve $C_{a,b,c}$ over an algebraically
closed field of characteristic two, there exists a connected cyclic
\'etale cubic cover with a supersingular elliptic quotient by a lift
of the hyperelliptic involution exactly when (C2.1) holds.
If (C2.1) fails, every connected cyclic \'etale cubic cover has
ordinary Jacobian and admits only constant maps to abelian varieties
of $2$-rank zero.

For the $2$-rank-one genus-two normal form
$y^2+y=a x^3+b x+c/x$, with $ac\ne0$, every elliptic quotient
of a cyclic cubic \'etale cover by a lifted hyperelliptic
involution is ordinary as well.

In particular take $a=b=1$ and $c=\alpha\in\F_4$ with
$\alpha^2+\alpha+1=0$. Then $(a+b+c)^3+abc=1+\alpha\ne0$.
Thus the construction using a single cyclic cubic cover and
$2$-rank-zero abelian targets cannot settle this curve. This is an
obstruction to that specified construction; covers of other odd
degrees remain available.
\end{claim}
\begin{proof}
We check that the normal form used above exhausts cyclic cubic
\'etale covers. Such covers over an algebraically closed field are
classified, up to choice of generator, by nonzero
$\eta\in J_C[3]$. Riemann--Roch supplies an effective divisor
$D_0$ of degree two with $\eta=[D_0-2P_\infty]$. It is unique:
a degree-two divisor with two independent sections belongs to the
hyperelliptic pencil and would give $\eta=0$.
It cannot contain $P_\infty$. Otherwise
$\eta=[S-P_\infty]$ and $3\eta=0$ would give a function with
its only pole of order three at the Weierstrass point $P_\infty$,
where three is a gap. In the displayed model this also follows from
$L(3P_\infty)=\langle1,x\rangle$.

Choose a function with divisor $3D_0-6P_\infty$.
The space $L(6P_\infty)$ has basis $1,x,x^2,x^3,Y$.
Its $Y$ coefficient cannot vanish: otherwise the divisor is invariant
under the hyperelliptic involution. An invariant effective divisor of
degree two is either a hyperelliptic fibre or a sum of Weierstrass
points, whose class relative to $2P_\infty$ is killed by two.
Neither can represent the nonzero three-torsion class $\eta$.
After scaling, the function is therefore $A+Y$, with
\[
 A=Lx^3+dx^2+ex+k,\qquad L\ne0.
\]
Let $u=x^2+u_1x+u_0$ be the monic polynomial whose zero divisor on
$C$ is $D_0+\iota D_0$. Taking norms and comparing the leading
coefficient gives
\[
                     A^2+hA+f=L^2u^3.                       \tag{C2.5}
\]
The required cubic cover is $z^3=(A+Y)/L$, after a harmless
constant rescaling over the algebraically closed field. Thus no
three-torsion class of smaller reduced degree has been omitted.

The lift $\tau:z\mapsto u/z$ of the hyperelliptic involution has
elliptic quotient. Here is a ramification verification in
characteristic two. The cubic cover has genus four. Over each of
the three branch points of $C\to\mathbf P^1$, the composite
$S_3$-cover has three points with distinct order-two inertia groups.
Exactly one is fixed by the specified $\tau$, and its different
exponent is two, as in the original Artin--Schreier cover. There
are no other fixed points. Riemann--Hurwitz for the quotient by
$\tau$ therefore reads $6=2(2g(E)-2)+6$, so $g(E)=1$.
Its invariant functions $x,w=z+u/z$ give the integral plane cubic
\[
 L\bigl(W^3+(X^2+u_1XZ+u_0Z^2)W\bigr)=X^2Z+XZ^2.            \tag{C2.6}
\]
It has the genus-one function field just computed, so its arithmetic
genus one forces it to be smooth. The Cech calculation in A.C2.1
makes its Hasse scalar $Lu_1$. It is supersingular exactly when
$u_1=0$.

Assume $u_1=0$. The coefficients of $x^5,x,x^3$ in (C2.5) give,
respectively,
\[
                  L=a,\qquad k=b,\qquad e=d+a+b+c.
\]
The coefficients of $x^4,x^2,1$ then imply
\[
 d^2+d+a=a^2u_0,\qquad
 u_0^2+u_0+1=((b+c)/a)^2,\qquad u_0^3=(b/a)^2.               \tag{C2.7}
\]
Choose $v$ with $v^2=u_0$. Since squaring is injective,
(C2.7) gives $v^3=b/a$ and $v^2+v+1=(b+c)/a$.
Consequently $(v+1)^3=c/a$, which implies (C2.1).
Conversely A.C2.1 constructs a connected cubic cover with
supersingular quotient whenever (C2.1) holds.

For the stronger assertion let $D\to C$ be any connected cyclic
cubic cover, let $\delta$ generate its deck group, and use the lift
$\tau$ above. Its Prym variety $P$ has dimension two, and
\[
                       J_D\sim J_C\times E^2.              \tag{C2.8}
\]
We justify the isogeny also in characteristic two. Choose an auxiliary
prime different from two and three and work on rational Tate modules.
The invariant part for $\delta$ is $VJ_C$; its complement is $VP$.
The hyperelliptic involution acts as $-1$ on $VJ_C$. Thus the image
of $VE$ under pullback lies in $VP$ and equals $(VP)^\tau$.
The two subspaces $(VP)^\tau$ and
$\delta(VP)^\tau$ have zero intersection: an element fixed by the
two corresponding reflections is fixed by $\delta$, whereas $VP$
has no $\delta$-invariants. Each has dimension two, and together
they span $VP$. The two pullback maps $E\to P$ consequently give
an isogeny $E^2\to P$. Pullback and norm for $D\to C$, whose
composition is $[3]$, give (C2.8).

If (C2.1) fails, every elliptic quotient (C2.6) is ordinary. The
base curve is ordinary, so (C2.8) makes $J_D$ ordinary. An abelian
quotient of an ordinary abelian variety is ordinary, by Poincare
reducibility and additivity of $2$-rank under products and isogenies.
A positive-dimensional abelian subvariety of a $2$-rank-zero variety
has $2$-rank zero. Hence a homomorphism from $J_D$ to such a variety
has zero image. The Albanese property then makes every map from
$D$ to such a variety constant. Substituting $(1,1,\alpha)$ gives
the asserted explicit example.

For the $2$-rank-one normal form put $Y=xy$, so $h=x$ and
$f=a x^5+b x^3+c x$. The same reduced-divisor argument and
ramification calculation give (C2.5)--(C2.6), with $h=x$.
The right side of the quotient cubic is then $XZ^2$.
The $x^5$ coefficient of the norm on the left is $a$, whereas
the right coefficient is $L^2u_1$. Thus $u_1\ne0$. The
coefficient of $XWZ$ in the quotient cubic remains $Lu_1$,
so it is ordinary. The base curve has $2$-rank one: Cartier
fixes $dx/x$ and kills $dx$, its regular differential basis.
\end{proof}

\begin{claim}{A.C2.3}{Proved}
Let $K$ be finitely generated over $\F_2$ and let $s\in K^\times$.
For the smooth proper genus-two curve
\[
                 C_s:\quad y^2+y=sx^3+x+\frac{s^5}{x},       \tag{C2.9}
\]
which has $2$-rank one, the geometrically prime-to-two point-section
map is injective. The assertion persists after finite extensions,
for finite-extension forms, and for hyperbolic open subcurves.

Equivalently this covers the $2$-rank-one models
$y^2+y=a x^3+b x+c/x$ with $abc\ne0$ and $a^5=b^{16}c$.
There is a connected cyclic \'etale cover $D\to C_s$ of degree
five with two jointly point-injective maps to a genus-two curve
of $2$-rank zero.
\end{claim}
\begin{proof}
The two Artin--Schreier poles have orders three and one, so the
genus is two. With $Y=xy$ the equation is
\[
 Y^2+xY=f,\qquad f=sx^5+x^3+s^5x.
\]
The regular differentials $dx/x,dx$ have Cartier images $dx/x,0$;
hence the $2$-rank is one. The point $P_\infty$ above infinity
again has pole orders two and five for $x,Y$.

Over a finite extension choose $d$ with $d^2+d=s^8$, and put
\[
 u=x^2+s^2,\qquad
 A=x^5+s x^4+x^3+x^2+d x+s^5.
\]
The direct identity
\[
                        A^2+xA+f=u^5                        \tag{C2.10}
\]
gives $F\iota(F)=u^5$ for $F=A+Y$.
At $x=s$, the hyperelliptic map is unramified and the two values
of $F$ differ by $s\ne0$. Exactly one point $R$ has $F=0$.
Since $u$ has order two there and $F$ has pole order ten at
infinity, its divisor is $10R-10P_\infty$. Consequently
$z^5=F$ defines an \'etale cyclic cover $D$ of degree five.
It is connected: a fifth root of $F$ would have its only pole of
order two at $P_\infty$, hence belong to $\langle1,x\rangle$,
which cannot produce the nonzero $Y$ coefficient. The genus of
$D$ is six.

The involution $\tau:z\mapsto u/z$ lifts the hyperelliptic
involution. Put $w=z+u/z$ and $X=x+s$, so $u=X^2$. The
fifth-power identity in characteristic two gives
\[
 w^5+X^2w^3+X^4w=x.
\]
Set $t=X^2+Xw+w^2$. Its left side is $w t^2$, so $x=w t^2$.
Substituting this into the definition of $t$ gives
\[
 (t^4+t^2+1)w^2+s w+(s^2+t)=0.
\]
Writing $H=t^4+t^2+1$ and $V=Hw/s$ exhibits the quotient
function field as that of
\[
 B_s:\quad V^2+V=\frac{(t^4+t^2+1)(t+s^2)}{s^2}.             \tag{C2.11}
\]
Indeed $\Omega(D)^\tau=\Omega(x,w)$, since $z$ satisfies
$z^2+wz+u=0$, and the displayed rational transformations identify
$\Omega(x,w)$ with $\Omega(t,V)$. The right side of (C2.11) has
one pole of order five. Thus $B_s$ is smooth proper of genus two,
and its regular differentials are $dt,t\,dt$. Cartier sends them
to $0,dt$, respectively. It is nilpotent, so $B_s$ has $2$-rank
zero and $J_{B_s}(\overline K)[2^\infty]=0$.

Let $r_0:D\to B_s$ be the quotient map, and let $\delta$ generate
the cyclic deck group of $D/C_s$. Define $r_1=r_0\circ\delta$.
These two maps jointly separate all geometric points by a group
action argument which also handles the ramification points.
A fibre of a quotient by $\langle\tau\rangle$ is a geometric
orbit. If distinct $P,Q$ had both images equal, then
\[
             Q=\tau P=\delta^{-1}\tau\delta P.
\]
The dihedral identity $\tau\delta\tau=\delta^{-1}$ would force
$\delta^2P=P$. This is impossible because a nonidentity deck
transformation of the \'etale cyclic cover $D/C_s$ fixes no point.

Compose $r_0,r_1$ with an Abel embedding of $B_s$ into its
Jacobian, choosing a base point after a finite extension if needed.
The resulting two maps have targets with no geometric
$2$-primary torsion points and are still jointly injective.
The section-lifting and Kummer argument of A.C2.1, or A.10,
therefore proves the asserted point-section injectivity. Finally,
$x'=b x$ transforms $y^2+y=a x^3+b x+c/x$ into
$y^2+y=(a/b^3)x'^3+x'+bc/x'$. The relation
$a^5=b^{16}c$ makes $bc=(a/b^3)^5$, proving the equivalent
parameter statement. The assertions about open subcurves and
finite extensions follow as in A.C2.1.
\end{proof}
